\chapter{运行与提交对照表}
\label{app:records}

\section{分类训练记录}
下列轮次来自日志与保存规则。任务一每组均包含五个完成的折，
表中以从 1 开始的轮次报告验证峰值。验证出现同一最高值时取首次。
任务二使用日志的索引值加 1；原始摘要中的 \code{epoch=11} 对应第 12 轮。

\begin{table}[H]
\centering\small
\caption{植物分类修复后各配置的轮次核对}
\begin{tabularx}{\linewidth}{@{}YrrP{44mm}r@{}}
\toprule
配置标记 & 计划 & 各折实际 & 各折首次验证峰值 & 写出权重 \\
\midrule
\code{fix_b0_288} & $4+46$ & 50 & 45, 37, 14, 15, 17 & 50 \\
\code{fix_cnx320} & $4+40$ & 44 & 33, 29, 27, 27, 19 & 44 \\
\code{fix_cnx320p} & $4+40$ & 44 & 24, 34, 16, 19, 26 & 44 \\
\code{fix_cnxs320m} & $4+40$ & 44 & 35, 25, 37, 26, 23 & 44 \\
\code{fix_cnx384m} & $4+40$ & 44 & 37, 34, 15, 16, 42 & 44 \\
\bottomrule
\end{tabularx}
\end{table}

末轮保存规则意味着训练曲线峰值用于诊断而不用于替换 OOF 权重。
比较上述配置时以五折实际 OOF 预测为单位。最终 Kaggle 的 v13 与 v16
在本地字节一致，故可用同一个 SHA-256 检查这两个文件。各自的提交编号仍不同。

\begin{table}[H]
\centering\small
\caption{任务二 v2 的模型组融合权重}
\begin{tabularx}{.88\linewidth}{@{}Yrr@{}}
\toprule
骨干 & 输入边长 & 权重 \\
\midrule
ConvNeXt-Small & 112 & $1/9$ \\
ConvNeXt-Small & 224 & $3/9$ \\
ConvNeXt-Tiny & 112 & $2/9$ \\
ResEmoteNet & 48 & $2/9$ \\
EfficientNetV2-S & 112 & $1/9$ \\
\bottomrule
\end{tabularx}
\end{table}

v2 权重取自对应的融合 JSON，类别先验指数为 0。单折的
\code{cns224_f0} 训练 48 轮并选取第 12 轮 EMA 权重。
这个权重只对应组内一折，表中的组权重施加于组级概率。

\clearpage
\section{定位与分割训练记录}
\begin{table}[H]
\centering\small
\caption{中心凹定位三个随机种子的实际与最优轮次}
\begin{tabularx}{\linewidth}{@{}lrrY@{}}
\toprule
种子 & 计划 & 各折实际 & 最优权重的 epoch（折 0 至 9） \\
\midrule
777  & 60 & 60 & 29, 41, 59, 45, 54, 20, 58, 12, 50, 26 \\
123  & 60 & 60 & 22, 14, 54, 59, 26, 10, 33, 21, 34, 46 \\
2026 & 60 & 60 & 33, 17, 14, 32, 14, 41, 41, 21, 31, 53 \\
\bottomrule
\end{tabularx}
\end{table}

三个随机种子的日志均覆盖 60 轮。按每折样本数 $n_f$ 和折内误差 $R_f$
恢复总体 OOF 的方法为
\[
 R_{\rm OOF}=\sqrt{\frac{\sum_{f=0}^9 n_f R_f^2}{\sum_{f=0}^9 n_f}}.
\]
由折外坐标得到种子 777、123、2026 的 RMSE 分别约为
14.17388、13.59294、15.02157 px，四舍五入后与历史汇总一致。
最后一轮日志可用于确认训练是否完成，复现推理时应使用
\code{best_f*.pt}。这些日志用于核查 B1 分支，VascX 直接调用
上游已训练权重。

\begin{table}[H]
\centering\small
\caption{血管分割中三条独立的权重记录}
\begin{tabularx}{\linewidth}{@{}P{35mm}P{32mm}P{30mm}Y@{}}
\toprule
运行 & 计划 & 实际记录 & 权重对应 \\
\midrule
FR-UNet 冻结方案 & 上游计划待核查 & 上游权重标识 40 & 课程提交使用本地封装的上游权重 \\
原生 Reti U-Net & 40 轮运行 & 日志 40 轮 & 第 15 轮原生验证 Dice 最高 \\
内部 Small U-Net & YAML 配置 120 轮 & 历史日志 11 轮 & 该运行最佳轮次记录为第 8 轮 \\
\bottomrule
\end{tabularx}
\end{table}

内部 Small U-Net 的 YAML 是当前留存配置。仅凭轮次数量不能确定
实际命令是否覆盖该参数，也不能判断结束原因。自训原生模型和内部
Small U-Net 的分数按各自日志和验证口径解释。

\clearpage
\section{主结果的提交编号与内容校验}
提交历史同时保存编号、文件名、时间和得分。表中列出原课程工作区
可直接核对的文件。归档数据表保留完整的 64 位十六进制 SHA-256；
本页仅展示前 12 位，便于查找。任务四的快照与原根目录
\code{submission.csv} 字节一致。历史提交中同名文件的其他得分
通过提交编号区分。

\begin{table}[H]
\centering\small
\caption{主结果对应的文件、Kaggle 编号及本地哈希前缀}
\begin{tabularx}{\linewidth}{@{}P{21mm}Yrr@{}}
\toprule
任务 & 结果文件 & 提交编号 & SHA-256 前缀 \\
\midrule
幼苗 & \code{submission_v16_4model.csv} & 56983147 & \texttt{8d1d42ed4305} \\
表情 & \code{sub_v2_cns224.csv} & 57308701 & \texttt{e85fbdf08896} \\
表情 & \code{sub_bagged.csv} & 57332301 & \texttt{792bbc68745e} \\
定位 & \code{submission_v15.csv} & 57701069 & \texttt{9f2352134700} \\
分割 & \code{submission_fr_official.csv} & 57946948 & \texttt{1b4415c5293e} \\
\bottomrule
\end{tabularx}
\end{table}

\subsection{最小格式校验}
\begin{description}[style=nextline,leftmargin=4.5em]
\item[分类] 逐行比较样例 ID 与实际提交 ID，检查 378 和 7\,065 行的
样本量，校验植物名称、表情整数标签和 ID 唯一性。
\item[定位] 逐行匹配 40 个坐标键，确认值为有限数且处于原图横纵边界。
使用 $(3,4)$ 与 $(0,0)$ 的示例核验实现输出 $\sqrt{12.5}$ 像素。
\item[分割] 读取 20 行，逐行解码为 $512\times512$ 掩膜并重新编码，
确认游程起点、长度、无重叠及往返一致性。最后校验冻结文件的 SHA-256。
\end{description}

这些校验直接读取已有的正式提交。Kaggle 分数源于平台保存的评估，
本地格式验证不涉及课程测试真值。
